% ============================================================
%  KHÁM PHÁ DỮ LIỆU (EDA)
% ============================================================
\section{Khám phá dữ liệu (Exploratory Data Analysis)}

% ------------------------------------------------------------
\subsection{Tổng quan dataset và phân phối biến mục tiêu}

OULAD gồm 7 bảng quan hệ, phủ 32.593 lượt đăng ký học phần với
10.655.280 lượt tương tác VLE. Bài toán phân loại 4 nhãn
\texttt{final\_result} $\in$ \{Pass, Withdrawn, Fail, Distinction\}.

\begin{center}
\small
\begin{tabular}{|l|r|l|}
\hline
\textbf{File} & \textbf{Dòng} & \textbf{Cột có null (\%)} \\
\hline
\texttt{studentInfo}            & 32.593      & \texttt{imd\_band} 3,4\% \\
\texttt{studentRegistration}    & 32.593      & \texttt{date\_unreg} 69,1\% \\
\texttt{studentAssessment}      & 173.912     & \texttt{score} 0,1\% \\
\texttt{studentVle}             & 10.655.280  & --- \\
\texttt{assessments}            & 206         & \texttt{date} 5,3\% \\
\texttt{courses}                & 22          & --- \\
\texttt{vle}                    & 6.364       & \texttt{week\_from/to} 82,4\% \\
\hline
\end{tabular}
\end{center}

\begin{center}
\small
\begin{tabular}{|l|r|r|}
\hline
\textbf{Nhãn} & \textbf{Số lượng} & \textbf{Tỉ lệ} \\
\hline
Pass        & 12.361 & 37,93\% \\
Withdrawn   & 10.156 & 31,16\% \\
Fail        &  7.052 & 21,64\% \\
Distinction &  3.024 & \phantom{0}9,28\% \\
\hline
\end{tabular}
\end{center}

Distinction chiếm $9,28\%$ — class nhỏ nhất, ít hơn Pass $\sim 4$ lần
$\Rightarrow$ chọn \textbf{Macro-F1} làm metric chính để không bỏ qua
class thiểu số.

% ------------------------------------------------------------
\subsection{Phân tích null ngữ nghĩa và kiểm toán rò rỉ dữ liệu}

\subsubsection*{2.1\enspace Semantic Null — null không phải thiếu data}

Sai lầm phổ biến là fillna/drop tất cả null cơ học. Trong OULAD mỗi
null đều mang ngữ nghĩa cụ thể:

\begin{center}
\small
\begin{tabular}{|l|r|p{3.6cm}|p{3.4cm}|}
\hline
\textbf{Cột} & \textbf{Null \%} & \textbf{Ý nghĩa thực} & \textbf{Quyết định} \\
\hline
\texttt{date\_unregistration}   & 69,1\% & NULL = học đến hết module  & \textbf{Drop} (leakage 99,1\% Withdrawn) \\
\texttt{score}                  & 0,1\%  & Bài không nộp              & \texttt{fillna(0)} \\
\texttt{imd\_band}              & 3,4\%  & Địa chỉ ngoài UK           & Encode \texttt{"Unknown"} \\
\texttt{week\_from / week\_to}  & 82,4\% & Resource trải toàn module  & Drop cả 2 cột \\
\hline
\end{tabular}
\end{center}

\texttt{date\_unregistration} có giá trị ở \textbf{99,1\%} ca Withdrawn
$\Rightarrow$ proxy gần tuyệt đối của label, bắt buộc drop để tránh
leakage 100\%. Ngược lại, NULL ở 69,1\% còn lại là tín hiệu
\textit{positive} (sinh viên tiếp tục học).

\subsubsection*{2.2\enspace Chọn cut-off theo data}

Median module length 262 ngày. Phân tích phân phối ngày rút môn trên
$n=10.156$ Withdrawn: day~135 ($51,6\%$ module) capture $\ge 80\%$
Withdrawn và còn $\sim 48\%$ thời gian can thiệp; day~175 capture
$\ge 90\%$ nhưng chỉ còn $33,1\%$ thời gian $\Rightarrow$ quá trễ.
Vì pipeline đánh giá trên nhiều mốc cảnh báo, chọn:
\textbf{cut-off giữa kỳ = 135} (mốc chính, tối ưu coverage vs khả
năng can thiệp); \textbf{cut-off cuối kỳ = \{30, 60, 90, 135\}} để
theo dõi cách tín hiệu tích luỹ theo thời gian khi đánh giá lại ở
cuối module.

\subsubsection*{2.3\enspace Kiểm toán leakage tại cut-off 135}

Toàn bộ 206 bài đánh giá phân loại SAFE (\texttt{deadline} $\le 135$)
vs LEAKAGE (\texttt{deadline} $> 135$):

\begin{center}
\small
\begin{tabular}{|l|r|r|r|p{3.4cm}|}
\hline
\textbf{Loại} & \textbf{SAFE} & \textbf{LEAK} & \textbf{Median submit ratio} & \textbf{Quyết định} \\
\hline
TMA  & \textbf{71} & 35 & 0,353 & Dùng (đa số SAFE) \\
CMA  & 17 & \textbf{59} & 0,577 & Chuẩn hóa theo từng module \\
Exam & 0  & \textbf{6}  & 0,927 & Drop hoàn toàn \\
\hline
\end{tabular}
\end{center}

Median submit ratio Exam ($0{,}927$) $\Rightarrow$ bài thi luôn ở cuối
kỳ, không có bài SAFE nào tại day~135, drop. CMA có 17/76 SAFE phân bố
không đều giữa modules (AAA, EEE không có CMA) $\Rightarrow$ FE chuẩn
hoá thành \texttt{cma\_submission\_rate}.

% ------------------------------------------------------------
\subsection{Phân tích tín hiệu feature}

\subsubsection*{3.1\enspace Tín hiệu nhân khẩu học}

Chi-square test ($n = 32.593$) trên các biến categorical so với nhãn
\texttt{final\_result}:

\begin{center}
\small
\begin{tabular}{|l|r|l|}
\hline
\textbf{Feature} & \textbf{$\chi^2$} & \textbf{Signal} \\
\hline
\texttt{code\_module}       & 1443,2 & \textbf{Rất mạnh} \\
\texttt{highest\_education} & 1024,7 & \textbf{Rất mạnh} \\
\texttt{imd\_band}          &  666,9 & \textbf{Rất mạnh} \\
\texttt{region}             &  449,7 & Mạnh \\
\texttt{age\_band}          &  222,7 & Mạnh \\
\texttt{disability}         &  138,5 & Trung bình \\
\texttt{gender}             &   16,5 & \textbf{Rất yếu} $\to$ drop \\
\hline
\end{tabular}
\end{center}

\texttt{gender} có $\chi^2$ thấp hơn \texttt{code\_module} gần
\textbf{90 lần} $\Rightarrow$ drop. 4 biến mạnh được giữ và encode
ordinal để bảo toàn quan hệ monotonic cho tree-based model.

\subsubsection*{3.2\enspace Hiệu ứng kỳ học (Semester B vs J)}

Pass rate kỳ B (tháng 2) $= 43{,}67\%$ ($n=12.488$), kỳ J (tháng 10)
$= 49{,}40\%$ ($n=20.105$). Chênh lệch không đồng đều theo module: CCC
diff $= -6{,}38$ pp, EEE diff $= -6{,}28$ pp (vượt ngưỡng $\pm 5$ pp)
$\Rightarrow$ semester \textbf{phải là feature}, không lọc bỏ.

% ------------------------------------------------------------
\subsection*{Kết luận EDA}

EDA xác định các nhóm tín hiệu khả dụng tại cut-off 135: (i)~nhân
khẩu học (4 biến sau chi-square), (ii)~điểm có trọng số trước
cut-off (TMA SAFE 71/106), (iii)~hành vi nộp bài (submission rate,
deadline gap). Tín hiệu hành vi VLE sẽ phân tích chi tiết ở giai
đoạn Feature Engineering sau khi thống nhất pipeline nhị phân
(At-Risk vs Not At-Risk).

\newpage
